\documentclass[10pt,letterpaper]{article}
\usepackage[T1]{fontenc}
\usepackage[utf8]{inputenc}
\usepackage{lmodern,microtype,amsmath,amssymb,booktabs,enumitem,xcolor,fancyhdr}
\usepackage[margin=1in,top=.88in,bottom=.88in,headheight=14pt]{geometry}
\definecolor{ink}{HTML}{17364C}
\definecolor{accent}{HTML}{006760}
\usepackage[colorlinks=true,urlcolor=accent,linkcolor=ink,pdftitle={Beyond Dirichlet GRH: CM coefficient fields and one-sided prime estimates}]{hyperref}
\usepackage{titlesec}
\titleformat{\section}{\large\sffamily\bfseries\color{ink}}{\thesection}{.65em}{}
\titleformat{\subsection}{\normalsize\sffamily\bfseries\color{ink}}{\thesubsection}{.65em}{}
\setlist{nosep,topsep=4pt}
\pagestyle{fancy}
\fancyhf{}\lhead{\small\sffamily Beyond Dirichlet GRH}\rhead{\small\sffamily Arithmetic research note}
\cfoot{\small\thepage}
\setlength{\parindent}{1.1em}
\setlength{\emergencystretch}{2em}
\newcommand{\repo}[2]{\href{https://github.com/lsernau/grh-lean/blob/72fead102f663d5c7a407964a13dc0cf9d977207/grh/#1}{#2}}
\begin{document}
\thispagestyle{plain}
{\sffamily\small\color{accent} RESEARCH NOTE\hfill 25 NOVEMBER 2026}\par\vspace{.8em}
{\sffamily\bfseries\LARGE\color{ink} Beyond Dirichlet GRH}\par\vspace{.4em}
{\Large Coefficient fields of CM newforms and one-sided prime estimates}\par\vspace{.9em}
\noindent\textbf{Main result.}
Exactly 26 quintic fields occur among the fields generated by Fourier coefficients
of weight-two CM newforms of all levels and all nebentypus. We give a defining
polynomial for every field and a finite certificate distinguishing every pair.
In degree seven the corresponding count is 31. For each prime degree
$\ell\ge7$, there is one field for each imaginary quadratic field of class number
$\ell$. These classifications are deductions in this note. They use the
algebraic theory of CM and unconditional class-number results. Martin explicitly
suggested degree five as a next case after his quadratic and cubic tables.
\href{https://arxiv.org/html/2509.24119v2}{Martin, Remark 1.2}.

\medskip\noindent
The starting point was the formal proof of Dirichlet GRH in
\href{https://github.com/lsernau/grh-lean/tree/72fead102f663d5c7a407964a13dc0cf9d977207}{\texttt{lsernau/grh-lean}}.
It supplies the quadratic-character hypothesis in Martin's finiteness argument.
In weight two an earlier theorem of Yang, applied to modular forms by
Gonz\'alez, gives a sharper restriction: the class number divides the
coefficient degree. Following that restriction leads to the prime-degree
classification. We also give a genus-field proof in even weight, derive
Grosswald's conjecture using the formal theorem, and adapt the repository's
one-sided prime-sum argument to a genuinely nonabelian test case.
The cubic-field test is an equivalence; its required prime estimate remains open
in this analysis. \href{https://arxiv.org/pdf/math/0301306}{Yang, Theorem 3.1};
\href{https://ems.press/content/serial-article-files/38349}{Gonz\'alez, Section 3}.
\medskip
\begin{center}\small
\begin{tabular}{p{.22\linewidth}p{.68\linewidth}}
\toprule
\textbf{Evidence} & \textbf{What it establishes}\\
\midrule
Lean source and independent replay & Five public RH and Dirichlet-GRH conclusions for the pinned repository revision.\\[3pt]
Published implications & Grosswald's bound and norm-Euclidean classifications follow by checking their analytic hypotheses.\\[3pt]
Proofs in this note & Prime-degree CM field classification, genus-field argument, weighted races, and a signed cubic splitting criterion.\\[3pt]
Exact computations & Polynomial irreducibility, distinctness of the 26 fields, class-number-five enumeration up to the published cutoff, and small primitive-root cases.\\
\bottomrule
\end{tabular}
\end{center}
\medskip\noindent
For the polynomial list and proofs see Sections~\ref{sec:quintic}--\ref{sec:oddprime}.
For the analytic transfer beyond the Dirichlet family see Section~\ref{sec:cubic}.

\medskip\noindent The arithmetic arguments downstream of the public GRH statement are given
as human-readable proofs. No end-to-end Lean formalization of those deductions
is claimed here. We found no earlier source tabulating the same 26 quintic
fields; the classification builds on results of Martin, Yang, Gonz\'alez
and the class-number literature, and we make no claim of publication priority.
\clearpage
\section{What the Lean repository establishes}\label{sec:audit}
The source revision is \texttt{72fead102f663d5c7a407964a13dc0cf9d977207}.
Both editions use Lean 4.35.0 and Mathlib
\texttt{56a005855ea17e0b0958f73ec801ae140b0a1679}.
The public statements import Mathlib alone. The two copies are byte-identical;
the five theorems in each edition's \texttt{Main.lean} inhabit those types
or Mathlib's own \texttt{RiemannHypothesis}. The four additional statements
give the zeta zero formulation excluding the negative even zeros and $s=1$,
Dirichlet half-plane nonvanishing, localization of all Dirichlet zeros in
the open critical strip, and its specialization to primitive characters.
Explicitly, for any modulus $N\ge1$ and character $\chi$ modulo $N$,
\[
 \Re s>\tfrac12,\quad (\chi\ne1\ \text{or}\ s\ne1)
   \quad\Longrightarrow\quad L(s,\chi)\ne0 .
\]
Here $\chi=1$ means the principal character. Imprimitive characters are included.
The public theorem says nothing about central multiplicities. For a finite
abelian extension of $\mathbb Q$, the factorization of its Dedekind zeta
function into primitive Dirichlet factors gives the corresponding
critical-strip result immediately.
\repo{src/Statements.lean}{Public definitions};
\repo{src/Main.lean}{five theorems};
\repo{scripts/PrintAxioms.lean}{axiom-report commands}.

\paragraph{Verification.}
We checked the revision, toolchain and dependency pins, theorem names,
and statement identity in a fresh checkout.
Andrew Sutherland records independent builds, replay of all 803 and 794 source
modules respectively, comparison against the challenge statements, and successful
multi-kernel checking of both editions. All five theorems have the standard
axioms \texttt{propext}, \texttt{Classical.choice}, and \texttt{Quot.sound}.
The repository's separate challenge template supplies the target types.
The checks cited here are for this shortened revision, not solely for the precursor.
We used Sutherland's recorded kernel results. We did not rerun the full Lean
build locally. \href{https://math.mit.edu/~drew/shortgrh/}{Sutherland, verification files for the shorter proof};
\repo{README.md}{project build instructions}.

\paragraph{How geometry yields a prime estimate.}
Let $q$ be a multiple of 18480 and let $d$ be a prime in the class
$q/2-1\pmod q$. The geometric object ultimately controls the remainder
\[
 R_{d,q}(x)=\sum_{\substack{n\le x\\n\equiv1\pmod q}}
  \Lambda(n)(1-(n/d))(1/n-1/x)-\frac{\log x}{\varphi(q)}-c_{d,q},
\]
with $(n/d)$ the Legendre symbol and the constant chosen so that the remainder
tends to zero. Theta derivatives put the weighted prime divisibilities into
a lattice of automorphic forms. Short integral forms give an upper bound for
a covolume corrected by an index. A spanning argument supplies sufficiently
many independent forms; in the final comparison at least half the target
dimension is represented. Determinants of chains of theta lifts, with
twists controlling the local denominators, give a lower bound for the same
quantity. Subtracting the common main term leaves
$m\,\mathrm{rem}\le f a(\log b)^5$ in the source notation.
The number of lifts satisfies $\sqrt x\le m\le2\sqrt x$.
The auxiliary scales and a fixed-modulus prime number theorem at a much larger
cutoff give $R_{d,q}(x)\le x^{-1/2+\epsilon}$ eventually for every
$\epsilon>0$. The square root in this estimate depends on the arithmetic
construction of the lifts.
\repo{GUIDE.md}{Reading guide};
\repo{src/New/Reduction/MainInequality.lean}{main comparison};
\repo{src/New/Reduction/OneSidedBound.lean}{one-sided bound}.

\paragraph{How one side detects a zero.}
The transform of this remainder has residues involving orders of untwisted
and quadratically twisted $L$-functions. Given a putative zero $\rho$, the
source chooses $d$ so none of the finitely many twisted characters has a zero
in the relevant rectangle to the right of $\beta>1/2$ and up to
$|\Im\rho|$. The remaining residue is proportional to
$\sum_\chi \operatorname{ord}_\rho L(s,\chi)/[\rho(1-\rho)]$, using
pole-corrected functions where appropriate. Real poles therefore have the
sign needed for Landau's lemma. An eventual upper bound then excludes the
putative pole, hence the zero. Taking modulus $18480N$ and changing levels
recovers any Dirichlet character; the extra Euler factors are harmless in
the relevant half-plane. Section~\ref{sec:cubic} reuses this last analytic
principle. The particular geometric prime estimate concerns a congruence class
and a quadratic character; replacing them by cubic splitting data is a new
arithmetic problem.
\repo{src/New/Reduction/Zeros/Criterion.lean}{Zero criterion};
\repo{src/Lib/Analytic/Landau.lean}{Landau lemma}.

\section{Coefficient fields of CM modular forms}
\label{sec:cm}
A normalized newform $f=\sum_{n\geq1}a_nq^n$ has a coefficient field
$K_f=\mathbb Q(a_1,a_2,\ldots)$. Fixing its weight and the degree of this
field still allows the level to grow without bound. Martin's Conjecture
1(a) asks whether only finitely many fields occur at fixed weight and
degree for trivial nebentypus (the Dirichlet character in the modular
transformation law). A family on which this can be decided consists of
the forms with complex multiplication, abbreviated CM. The decisive
zeros belong to quadratic Dirichlet $L$-functions. One needs no
hypothesis on the zeros of the modular forms' own $L$-functions.

\paragraph{Finiteness with varying character.}
Take the repository's Dirichlet GRH theorem as input. For fixed integers
$k\geq2$ and $d\geq1$, only finitely many isomorphism classes of fields
$K_f$ occur among CM newforms of weight $k$ with
$[K_f:\mathbb Q]=d$, even as both the level and nebentypus vary.

Martin proves the trivial-character case under quadratic Dirichlet GRH
(Theorem~1). His Remark~1.1 anticipates a fixed-character analogue and
points to the same methods for CM forms. We give the simultaneous
deduction, including the field identity that allows the character to
vary. The class-group and value-field arguments are Martin's.
\href{https://arxiv.org/pdf/2509.24119v2}{Martin, \emph{Rationality fields of CM modular forms}, Remark 1.1 and Sections 2--3}.

\paragraph{Recover the Hecke character's field.}
Write $E$ for the imaginary quadratic CM field and $\ell=k-1>0$.
The CM correspondence realizes $f$ as the theta series of an algebraic
Hecke character $\psi$ of $E$ of type $(\ell,0)$ and conductor
$\mathfrak m$. Thus $a_n$ is the sum of $\psi(\mathfrak a)$ over integral
ideals of norm $n$, extending $\psi$ by zero at ideals meeting the
conductor. Let $L$ be the field generated by its values and $K=K_f$.
Then $K\subset L$. We claim that $L=EK$ for every nebentypus.

For $p=\mathfrak p\bar{\mathfrak p}$ split and prime to the conductor,
put $u=\psi(\mathfrak p)$ and $v=\psi(\bar{\mathfrak p})$. The
coefficients at $p$ and $p^2$ give
\[
       u+v=a_p,\qquad uv=a_p^2-a_{p^2}.
\]
Take $h$ to be the exponent of the ray class group modulo $\mathfrak m$.
Then $\mathfrak p^h=(\alpha)$ and
$\bar{\mathfrak p}^{h}=(\beta)$ with generators congruent to $1$
modulo $\mathfrak m$. Hence $u^h=\alpha^\ell$ and $v^h=\beta^\ell$
belong to $E$. Their principal ideals are distinct, and neither is
invariant under conjugation. Consequently $u^h$ generates $E$ over
$\mathbb Q$, so $E\subset L$. Every automorphism of
$\overline{\mathbb Q}$ fixing $EK$ preserves the pair of roots $u,v$ of
$X^2-a_p X+(a_p^2-a_{p^2})$ and fixes their distinct $h$th powers.
It therefore fixes each root. This gives $u,v\in EK$. At a split prime
where only one factor is prime to $\mathfrak m$, its value equals $a_p$.
At an inert prime the sole prime ideal has value $a_{p^2}$, and at a
ramified prime its value is $a_p$, whenever it is prime to the conductor.
Prime ideals generate the ideal group, which proves the claim. In
particular $[L:E]=d/[K\cap E:\mathbb Q]$ divides $d$.

\paragraph{Bound the CM discriminant.}
For an ideal $\mathfrak a$ prime to $\mathfrak m$, write
$\mathfrak a^h=(\gamma)$ with $\gamma\equiv1\pmod{\mathfrak m}$.
Taking $h$th powers of ideals and using $\psi(\mathfrak a)^h=\gamma^\ell$
gives
$(\psi(\mathfrak a))=(\mathfrak a\mathcal O_L)^\ell$ in $L$.
The ideal norm from $L$ to $E$ shows that
$\mathfrak a^{\ell[L:E]}$ is principal. Since every class has such a
representative, $e_E=\exp\mathrm{Cl}(E)$ divides $(k-1)d$.

Here is where the analytic result enters. Write $D$ for the fundamental
discriminant of $E$. If $p=\mathfrak p\bar{\mathfrak p}$ splits, take
$\mathfrak p^{e_E}=(\alpha)$. Unequal valuations at the two primes above
$p$ force $\alpha\notin\mathbb Q$. Writing
$\alpha=(a+b\sqrt D)/2$ with $b\ne0$ gives
\[
       p^{e_E}=|N_{E/\mathbb Q}\alpha|\geq |D|/4 .
\]
A sufficiently small split prime therefore bounds the discriminant of every possible $E$.
A numerical version gives a test for an individual proposed CM field. Put
$T=\log|D|$ and $M=(k-1)d$. If $T>25$, Theorem~3 of Elsenhans,
Kl\"uners and Nicolae gives $p\leq(1.881T+6.18)^2$ under ERH for the
quadratic field. Combining it with the preceding ideal-norm inequality
forces
\[
               T\leq\log 4+2M\log(1.881T+6.18).
\]
The right side grows only logarithmically with $T$, so this is an
explicit discriminant cutoff for every weight and coefficient degree.
For example, in weight two and degree seven the class-group exponent
divides seven, and their Table~2 yields $|D|\leq 3.9\cdot10^{30}$.
\href{https://arxiv.org/pdf/1803.02056}{Elsenhans, Kl\"uners and Nicolae, Theorem~3 and Table~2}.

\paragraph{Finish the field count.}
For any one of these $E$, Martin's Corollary~3.6 gives finitely many
value fields $L/E$ of each bounded degree, uniformly in the conductor.
Indeed the image on principal ideals generates $E(\mu_r)$, where $r$
is the order of the finite modulus character. Thus
$\varphi(r)\leq [L:\mathbb Q]\leq 2d$. Choose ideal-class generators
$\mathfrak t_i$ with $\mathfrak t_i^{n_i}=(\theta_i)$.
Martin's Proposition~3.4 places $L$ in the field obtained by adjoining
$\mu_{r e_E}$ and the $n_i$th roots of $\theta_i^\ell$ to $E$, independently
of $\mathfrak m$. There are finitely many possible $r$, and each resulting
number field has only finitely many subfields. This proves the assertion.

\paragraph{Why weight two admits more.}
For trivial nebentypus Martin's Theorem~2 lists 23 real quadratic and 18 cubic
fields using quadratic Dirichlet GRH. In weight two the divisibility
$h_E\mid[L:E]\mid d$ of Yang and Gonz\'alez gives an independent, unconditional
reduction from class-group exponent to class number. Gonz\'alez's Theorem~3.1
already proves weight-two finiteness in all degrees and for all characters.
Our interest here is the exact list over all levels. The next sections give the
genus-field and cubic arguments, then the classification in higher prime degrees.
For general weights the split-prime proof above works whenever a
sufficiently small split prime is available. The full Dirichlet GRH supplies
one route. Martin's conjecture for non-CM newforms requires other ideas.
\href{https://arxiv.org/pdf/2509.24119v2}{Martin, Theorems 1--2};
\href{https://ems.press/content/serial-article-files/38349}{Gonz\'alez, Theorems 3.1--3.2}.

\section{The genus field inside the value field of a CM newform}
\label{sec:genus}

After Theorem~2 of his paper, Martin observes that when a weight-two CM newform of trivial nebentypus has quadratic coefficient field $K$ and its imaginary quadratic CM field $E$ has class number two, $EK$ is the Hilbert class field of $E$. He obtains the equality from the field tables and singles out the lack of a non-computational proof. The following containment explains the equality, extends to every even weight, and removes the analytic hypothesis from his classification in coefficient degree two. The ramified-prime square-root identity in the proof has a weight-two antecedent in Yang's Lemma~3.3. Gonz\'alez's earlier treatment of the 83 quadratic pairs already identifies the weight-two fields, including the class-number-two case. The following argument provides an alternative proof. Its local sign uses the trivial nebentypus and works in every even weight. See \href{https://arxiv.org/pdf/math/0301306}{Yang, \emph{On CM abelian varieties over imaginary quadratic fields}}. See \href{https://arxiv.org/html/2509.24119v2}{Martin, introduction, Lemma~2.2 and Sections~5--7}.

\paragraph{Theorem.} Let $f$ be a normalized CM newform of trivial nebentypus and even weight $k\ge2$, with imaginary quadratic CM field $E$, coefficient field $K$, and Hecke-character value field $L$. Write $D<0$ for the fundamental discriminant of $E$ and $t$ for the number of distinct primes dividing $D$. The genus field $G_E$ of $E$ satisfies
\[
        G_E\subseteq L=EK,\qquad G_E\cap\mathbb R\subseteq K,
        \qquad 2^{t-1}\mid[K:\mathbb Q].
\]
In particular, if $h_E=[K:\mathbb Q]=2$, then $EK=H_E$, the Hilbert class field, and $K$ has discriminant equal to the positive signed prime-discriminant factor of $D$.

\paragraph{The sign at an odd ramified prime.}
Write $w=k-1$, which is odd. The algebraic Hecke character $\psi$ attached to $f$ has infinity type $(w,0)$ and conductor $\mathfrak f$. Its finite character $\eta:(\mathcal O_E/\mathfrak f)^\times\to\mathbb C^\times$ satisfies
\[
       \psi((a))=\eta(a)a^w\quad (a\in E^\times\text{ a unit at }\mathfrak f).
\]
The restriction of $\eta$ to rational units is the primitive quadratic character $\chi_D$ away from the moduli, by the formula for the nebentypus of $f_\psi$. Its finite image $U$ lies in $L$ by the displayed identity, since $E\subset L$ and each residue unit has a lift. Group the local components of $\eta$ over each rational prime $v$ into $\eta_v$. Their restrictions to diagonal rational units are exactly the corresponding local components of $\chi_D$. This follows by factoring the equality of Dirichlet characters into prime-power components using the Chinese remainder theorem. The assertion allows $\mathfrak f$ to be asymmetric over a split prime.

Fix any odd prime $s\mid D$ and let $P_s^2=(s)$ be its ramified prime ideal. By weak approximation choose $x\in E^\times$ such that $\mathfrak a=xP_s$ is prime to $\mathfrak f$. Then $\mathfrak a^2=(\theta)$ for $\theta=sx^2$. The defining formula for $\eta$ gives
\begin{equation}
      b^2=s\eta(\theta),\qquad
      b=\frac{\psi(\mathfrak a)}{x^w s^{(w-1)/2}}\in L.   \label{eq:genus-square}
\end{equation}
We now compute $\eta(\theta)$ in $U/U^2$. At a prime $v\ne s$ in the conductor, $\theta$ equals $s$ times a local unit square at each relevant place. At $s$ use instead
\[
                 \theta=(D/s)(s x/\sqrt D)^2.
\]
Here $D/s$ is a rational local unit and the second factor is a local unit square. Put $e_s=(-1)^{(s-1)/2}$, $s^*=e_s s$, and $d'=D/s^*$. The factorization into primitive quadratic characters gives $\chi_D=\chi_{s^*}\chi_{d'}$ on common units, with coprime conductors. Away from $s$ the rational local components evaluated at $s$ multiply to $\chi_{d'}(s)=(\frac{d'}s)$. At $s$ the restriction to rational units is $(\frac{\cdot}s)$. Multiplying in $U/U^2$ yields
\[
     \eta(\theta)\equiv
        \left(\frac{d'}s\right)\left(\frac{D/s}{s}\right)
        =\left(\frac{e_s(d')^2}{s}\right)=e_s
        \pmod {U^2}.
\]
All signs belong to $U$ since $\eta(-1)=-1$. Consequently there is a $u\in U\subset L$ with $\eta(\theta)=e_s u^2$. Equation~\eqref{eq:genus-square} gives $(b/u)^2=s^*$. We have proved $\sqrt{s^*}\in L$ for every odd prime $s\mid D$, without any restrictions on the conductor, the order of $\eta$, or the class of $P_s$.

\paragraph{Recovering the genus field.}
Gauss genus theory gives $G_E=\mathbb Q(\sqrt{d_1},\ldots,\sqrt{d_t})$, where $D=\prod_j d_j$ is the factorization into signed prime discriminants, and $[G_E:E]=2^{t-1}$. Equivalently $\operatorname{Gal}(G_E/E)$ is the class group modulo squares. The odd factors are the $s^*$ just obtained. If a dyadic factor $-4,8,$ or $-8$ occurs, its square root also belongs to $L$ by dividing $\sqrt D\in E$ by the product of the odd-factor square roots. This proves $G_E\subseteq L$. Martin's Lemma~2.2 identifies $L=EK$, and $K$ is totally real. Thus $L\cap\mathbb R=K$ and $[L:E]=[K:\mathbb Q]$, giving both remaining conclusions. For $h_E=2$ the class group has trivial subgroup of squares, so $G_E=H_E$ has degree two over $E$. When $[K:\mathbb Q]=2$, containment is therefore equality. The two signed prime discriminants have opposite signs; the positive one defines the unique real quadratic subfield. For genus theory see \href{https://www.math.purdue.edu/~jlipman/MA598/Lemmermeyer-principal\%20genus.pdf}{Lemmermeyer, \emph{The development of the principal genus theorem}}. \hfill$\square$

\paragraph{Corollary in weight two.}
Suppose $k=2$ and $[K:\mathbb Q]=2$. Martin's Corollary~3.2 implies $\exp\operatorname{Cl}(E)\mid 2$. The class group hence equals its quotient by squares. By the theorem its order $2^{t-1}$ divides two. Thus $h_E$ is one or two. The unconditional classifications of imaginary quadratic fields of class numbers one and two provide exactly the possible base fields to check. Martin's local arguments and constructions for these fields now give the degree-two part of his Theorem~2 without ERH. Equivalently, the complete list of $K$ consists of the real quadratic fields with these 23 discriminants:
\[
\begin{gathered}
 5,8,12,13,17,21,24,28,29,33,37,41,\qquad\\[-2pt]
 44,57,61,76,89,129,172,201,268,489,652 .
\end{gathered}
\]
The list and its existence proofs are Martin's Table~1.1 and Sections~5--7. The genus argument excludes any further exponent-two fields by their class-group rank. Earlier work gives another route in weight two: Yang's Theorem~3.1 implies $h_E\mid[K:\mathbb Q]$, and Gonz\'alez's Theorem~3.2 already classifies the 83 pairs of CM field and quadratic coefficient field when nebentypus is allowed to vary. These results predate Martin's list; we claim no priority for the weight-two quadratics. See \href{https://ems.press/content/serial-article-files/38349}{Gonz\'alez, \emph{Endomorphism algebras of CM modular abelian varieties}}. The finite positive constructions are independent of ERH, as are the specific class-number-one exclusions. This deduction concerns trivial nebentypus and degree two; Martin's cubic classification has a different exponent obstruction. For the class-number-two classification see \href{https://magma.maths.usyd.edu.au/~watkins/papers/smallcn.pdf}{Watkins, \emph{Class numbers of imaginary quadratic fields}}.

\section{The cubic classification without a zero hypothesis}
\label{sec:class-three}

In weight two precisely eighteen cubic fields occur as coefficient
fields of CM newforms, across all levels and nebentypus. The
discriminants are
\[
\begin{gathered}
49,\ 81,\ 321,\ 621,\ 837,\ 993,\ 1593,\ 1929,\ 2241,\ 3753,\\
5697,\ 7641,\ 8289,\ 10233,\ 13473,\ 14769,\ 23841,\ 24489 .
\end{gathered}
\]
Martin's Table 1.2 identifies each field up to isomorphism by its
discriminant. His proof used a GRH-conditional list of fields whose
class-group exponent is three. The degree of the character's value
field lets us use the unconditional class-number-three list instead.

\paragraph{The algebraic restriction.}
A cubic field $K$ has a real embedding. Its embedding conjugates the
newform to one with real coefficients, so Sch\"utt's Corollary~1.2(i)
forces its nebentypus, and therefore the original one, to be trivial.
Write $E$ for the imaginary quadratic CM field, $L=EK$ for the
value field of the inducing character $\psi$, and $\eta$ for its finite
modulus character. Then $[L:E]=3$ and $\exp\mathrm{Cl}(E)\mid3$ by
the ideal-norm argument of Section~\ref{sec:cm}.
\href{https://arxiv.org/pdf/math/0511228}{Sch\"utt, Corollary~1.2}.

Assume $h_E>1$. We first show that $\eta$ has order two. If its order
is $r$, then $\mathbb Q(\zeta_r)\subseteq L$, so $\varphi(r)\mid6$.
Also $r$ is even since $\eta(-1)=-1$. If $\varphi(r)=2$, the cyclotomic
quadratic field would have to equal $E$, as $[L:E]$ is odd, but both
$\mathbb Q(i)$ and $\mathbb Q(\sqrt{-3})$ have class number one.
If $\varphi(r)=6$ then $r=14$ or $18$. The sextic cyclotomic field
would equal $L$, whose quadratic subfield would respectively be
$\mathbb Q(\sqrt{-7})$ or $\mathbb Q(\sqrt{-3})$, again of class
number one. Only $r=2$ remains.

For any ideal class $c$ choose a fractional ideal $\mathfrak a$
representing $c$ prime to the conductor and choose $\theta_c\in E^\times$
with $\mathfrak a^3=(\theta_c)$. The assignment
\[
\mathrm{Cl}(E)\longrightarrow E^\times/(E^\times)^3,\qquad
c\longmapsto [\theta_c]
\]
is a well-defined injective homomorphism. A change of representative
multiplies $\theta_c$ by a cube times a unit, and both units
$\pm1$ are cubes. The law follows from ideal multiplication.
If $\theta_c=b^3$, unique factorization of fractional ideals gives
$\mathfrak a=(b)$. This proves injectivity.

Every $\theta_c$ acquires a cube root in $L$ because
$\psi(\mathfrak a)^3=\eta(\theta_c)\theta_c=\pm\theta_c$.
Set $F=E(\zeta_3)$. Since $[F:E]=2$, the norm from $F$ to $E$ proves
that $E^\times/(E^\times)^3\longrightarrow F^\times/(F^\times)^3$
is injective: if $b$ becomes a cube, then $b^2$ is a cube in $E$,
so $b$ is a cube already. By the Kummer correspondence over $F$,
adjoining cube roots of a subgroup of $F^\times/(F^\times)^3$
of order $3^m$ gives an extension of degree $3^m$.
The subgroup just constructed has order $h_E$ and its radicals are
all in $FL$. It follows that
\[
               h_E\mid[FL:F]\mid3 .
\]
In particular every possible nontrivial class group here is cyclic of order
three. This excludes fields with arbitrarily many independent classes
of order three without enumerating them.

\paragraph{Finish from finite tables.}
There are sixteen imaginary quadratic fields of class number three,
with the largest absolute discriminant $907$. The unconditional
classification goes back to Oesterl\'e and is included in Watkins's
classification up through class number $100$.
For these sixteen fields Martin's Corollary~4.5 and Section~6.2
give the unique cubic coefficient field in each case, and
Table~6.3 gives its discriminant; the constructions there establish
occurrence. For the class-number-one case his Section~5.3 gives just
$\mathbb Q(\zeta_7)^+$ and $\mathbb Q(\zeta_9)^+$. These are precisely
the eighteen entries displayed above. The computation of the individual
fields is Martin's. Yang's divisibility, as applied by Gonz\'alez in Section~3 of his paper,
also gives $h_E\mid3$ directly. The Kummer argument above supplies
a separate proof in this degree. The change relative to Martin's proof is the uniform exclusion
of class groups of order greater than three.
\href{https://arxiv.org/pdf/2509.24119v2}{Martin, Sections~5.3 and~6.2, Tables~1.2 and~6.3};
\href{https://magma.maths.usyd.edu.au/~watkins/papers/smallcn.pdf}{Watkins, \emph{Class numbers of imaginary quadratic fields}}.

\section{A classification in degree five}
\label{sec:quintic}

\textbf{Theorem.} Across all levels and nebentypus, exactly twenty-six
quintic fields occur as coefficient fields of weight-two CM newforms.
One is the maximal real subfield of $\mathbb Q(\zeta_{11})$. Each of
the other twenty-five comes from a different imaginary quadratic field
of class number five. This result is unconditional. Martin singled out
degree five in Remark 1.2 as a possible extension of his explicit
classification. \href{https://arxiv.org/pdf/2509.24119v2}{Martin, Remark 1.2}.

\paragraph{The input from degree five.}
As in the cubic argument, the odd coefficient degree implies trivial
nebentypus. Put $L=EK_f$, so $[L:E]=5$ and $\exp\mathrm{Cl}(E)\mid5$.
For $h_E>1$, $E$ has only the units $\pm1$.
Let $r$ be the order of the finite character $\eta$.
As $\varphi(r)\mid10$ and $r$ is even, $r$ is one of $2,4,6,22$.
Orders four and six would force $E$ to be the quadratic cyclotomic
field inside $L$, since $[L:E]$ is odd. Order twenty-two would
force $L=\mathbb Q(\zeta_{11})$ and hence $E=\mathbb Q(\sqrt{-11})$.
All three exceptional quadratic fields have class number one.
Thus $r=2$.

For each class $[\mathfrak a]$, use a representative prime to the
conductor, and write $\mathfrak a^5=(\theta)$. The classes $[\theta]$
form a subgroup of $E^\times/(E^\times)^5$ isomorphic to
$\mathrm{Cl}(E)$ by the ideal argument of the preceding section.
Each $\theta$ has a fifth root in $L$ since $\psi(\mathfrak a)^5=\pm\theta$.
These classes stay independent over $F=E(\zeta_5)$, as $[F:E]$ divides
four. Kummer theory gives $h_E\mid[FL:F]=5$. The only cases to
consider are therefore $h_E=1$ and $h_E=5$.

\paragraph{Class number one.}
Principal ideals supply all values, so $L=E(\mu_r)$.
For this field to have degree ten over $\mathbb Q$ the preceding
cyclotomic argument leaves only $E=\mathbb Q(\sqrt{-11})$ and
$r=22$. To realize it, reduce modulo $11\mathcal O_E$. The quotient
is the ring of dual numbers over $\mathbb F_{11}$, with unit group
$\mathbb F_{11}^\times\times(1+\epsilon\mathbb F_{11})$,
$\epsilon^2=0$. Choose the Legendre character on the first factor
and a character of order eleven on the second. The resulting
$\eta$ sends $-1$ to $-1$ and restricts on rational units to
$\chi_{-11}$. Hence $\psi((\alpha))=\alpha\eta(\alpha)$ defines a
type-one Hecke character, primitive because the order-eleven part
is nontrivial. Its newform has trivial nebentypus and the asserted
real quintic coefficient field.

\paragraph{Class number five.}
Genus theory and the odd class number give a single prime discriminant,
necessarily $D=-q$ with $q\equiv3\pmod4$ an odd prime; the purely
dyadic cases have class number one. Let $\mathfrak t$ be an integral
ideal representing a generator, prime to $q$, with norm $n$.
Choose $\theta$ with $\mathfrak t^5=(\theta)$.
There is exactly one possible coefficient field up to isomorphism.
Indeed for any conductor take a representative $\alpha\mathfrak t$
coprime to it. Its character value, divided by $\alpha$, is a fifth
root of $\pm\theta$. The principal-ideal values generate $E$, and this
one value generates the rest of $L$ over $E$. The sign can be absorbed
in the choice of fifth root. The polynomial $X^5-\theta$ is
irreducible over $E$: if a root generated an extension of degree
$1\leq j<5$, its norm would show $\theta^j$ is a fifth power, forcing
$\theta$ itself to be a fifth power. This contradicts the order of
$\mathfrak t$. Thus all choices yield isomorphic $L$ over $E$.
As $L$ is a totally imaginary quadratic extension of its totally real
coefficient field, its maximal totally real subfield recovers $K_f$.

For existence, choose conductor the prime above $q$ and $\eta$ the
Legendre character on its residue field. It sends $-1$ to $-1$ and its
rational restriction is $\chi_D$. The principal-ideal rule
$\psi((\alpha))=\eta(\alpha)\alpha$ extends to all ideals prime to the
conductor by assigning to $\mathfrak t$ any fifth root of
$\eta(\theta)\theta$: every ideal is a principal ideal times a power
of $\mathfrak t$, and its only generator relation modulo principal
ideals is the fifth-power relation used here. This character is primitive
at $q$. Its newform has weight two, trivial nebentypus and coefficient
degree five. If $z=\psi(\mathfrak t)$ and $\Theta=z^5$, then
$\overline z=n/z$ because $|\theta|^2=n^5$. In particular the field
is generated by $x=z+n/z$, a root of
\[
X^5-5nX^3+5n^2X-\operatorname{Tr}_{E/\mathbb Q}(\Theta).
\]
Indeed $x$ is real and $[E(z):E]=5$, while $z$ satisfies
$Z^2-xZ+n=0$, so $x\notin\mathbb Q$.

\paragraph{Why different CM fields give different quintic fields.}
Set $B=E(\zeta_5)$ and $M=B(z)$. Since $[B:E]$ divides four,
$\Theta$ remains a non-fifth-power in $B$. Thus $M/B$ is cyclic of
degree five. Conjugation of $\Theta$ takes it to $n^5/\Theta$, so $M$
is normal over $\mathbb Q$. It is unramified outside $5q$: at all
other places $B$ is unramified and the valuation of $\Theta$ is a
multiple of five. Locally divide by a fifth power of a uniformizer
and adjoin a fifth root of a unit. Its reduction is separable at that
residue characteristic, giving an unramified extension.

At $q$, an inertia element restricting nontrivially to $E$ fixes
$\zeta_5$ and sends $z$ to $\zeta_5^j n/z$ for some $j$. The five
conjugates of $x$ are $x_i=\zeta_5^i z+\zeta_5^{-i}n/z$.
They are distinct, since an equality for two different $i$ would make
$z^2/n$ a root of unity, contrary to the different ideals
$\mathfrak t^5$ and $\overline{\mathfrak t}^{5}$.
Indeed equality of these ideals would imply $[\mathfrak t]=[\mathfrak
t]^{-1}$. The inertia element acts on the five conjugates by
$i\mapsto-i-j$ and therefore acts nontrivially. This forces
ramification in the quintic field: if it were unramified at $q$, the
compositum of all its local conjugates would be unramified there as
well. Thus $q$ is the unique ramified rational prime other than five
and recovers the CM field. The cyclotomic example ramifies only at
eleven, and $\mathbb Q(\sqrt{-11})$ has class number one, so it is
distinct from every class-number-five example.

The unconditional class-number-five classification contains twenty-five
fields, with largest absolute discriminant $2683$. This proves the
count. \href{https://magma.maths.usyd.edu.au/~watkins/papers/smallcn.pdf}{Watkins, \emph{Class numbers of imaginary quadratic fields}, final table}.

The canonical form of level $q^2$ and coefficient degree $h(-q)$ is
part of earlier work, including Gonz\'alez--Lario, Proposition~5.3 and Gross, Section~8.
Yang's class-number divisibility gives the reduction to class number one or
five directly. The additional assertions here are exhaustiveness over all
levels, distinction by the ramified prime, and the explicit polynomial list.
\href{https://arxiv.org/pdf/0810.4224}{Gonz\'alez and Lario, Proposition~5.3};
\href{https://people.math.harvard.edu/~gross/preprints/hecke.pdf}{Gross, Section~8}.

\subsection*{Defining polynomials and exact checks}
The class-number-five rows below use
\[
             P_{p,a}(X)=X^5-5pX^3+5p^2X-a .
\]
Here $D_E=-q$, $p$ is the norm of a nonprincipal split prime ideal, and
$4p^5=a^2+qb^2$. The integer $(a+b\sqrt{-q})/2$ generates a fifth
power of one of the two prime ideals over $p$. Indeed it has the
required norm and $p\nmid a$, so its two conjugates cannot both be divisible by the
same prime over $p$. Changing the sign of the trace sends the root to its
negative and gives the same field up to isomorphism.
\begin{center}
\begin{tabular}{rrrr@{\hspace{2.5em}}rrrr}
\toprule
$q$ & $p$ & $a$ & $b$ & $q$ & $p$ & $a$ & $b$\\
\midrule
47 & 2 & 9 & 1 & 691 & 5 & 38 & 4 \\
79 & 2 & 7 & 1 & 739 & 5 & 26 & 4 \\
103 & 2 & 5 & 1 & 787 & 7 & 59 & 9 \\
127 & 2 & 1 & 1 & 947 & 3 & 5 & 1 \\
131 & 3 & 29 & 1 & 1051 & 5 & 107 & 1 \\
179 & 3 & 16 & 2 & 1123 & 7 & 239 & 3 \\
227 & 3 & 8 & 2 & 1723 & 11 & 149 & 19 \\
347 & 3 & 25 & 1 & 1747 & 17 & 1565 & 43 \\
443 & 3 & 23 & 1 & 1867 & 7 & 4 & 6 \\
523 & 7 & 220 & 6 & 2203 & 7 & 255 & 1 \\
571 & 5 & 58 & 4 & 2347 & 17 & 1810 & 32 \\
619 & 5 & 109 & 1 & 2683 & 11 & 344 & 14 \\
683 & 3 & 17 & 1 &  &  &  &  \\
\bottomrule
\end{tabular}
\end{center}
The remaining field has defining polynomial
$X^5+X^4-4X^3-3X^2+3X+1$, obtained by writing
$1+\zeta+\cdots+\zeta^{10}=0$ in terms of $X=\zeta+\zeta^{-1}$ for
a primitive eleventh root. The first row of the table, for example,
gives $X^5-10X^3+20X-9$ for $D_E=-47$.

The program \texttt{checks/quintic\_fields.py} counts reduced primitive
binary quadratic forms through Watkins's discriminant bound, finds the
displayed norm identities, and checks every polynomial modulo primes.
For each of the 26 polynomials the archive gives a prime of irreducible
reduction. For each of the $\binom{26}{2}=325$ pairs it gives a prime
where both reductions are squarefree and the factor degrees disagree.
A squarefree reduction ensures the prime does not divide the polynomial
discriminant, hence does not divide the order index. Its factor degrees
therefore give the residue degrees in the field; unequal patterns certify
nonisomorphism. These checks verify the list and its distinctness.
The proof in Section~\ref{sec:quintic} establishes completeness over all levels.

\section{A classification in every odd prime degree at least seven}
\label{sec:oddprime}

\paragraph{Theorem.}
For each prime $\ell\ge7$, the isomorphism classes of degree-$\ell$ coefficient fields of weight-two CM newforms of trivial nebentypus are in bijection with the imaginary quadratic fields $E$ of class number $\ell$. By the real-embedding argument of Section~\ref{sec:class-three}, this includes every nebentypus. The map assigns to $E$ the real field generated by a root of the Dickson polynomial described below. In particular the number of such fields is $31$ in degree seven, $41$ in degree eleven and $37$ in degree thirteen. These counts use the unconditional class-number enumeration in \href{https://magma.maths.usyd.edu.au/~watkins/papers/smallcn.pdf}{Watkins, final class-number table}.

\paragraph{Restricting the finite character.}
Let $L=EK$ be the Hecke-character value field. Yang's class-number divisibility gives $h_E\mid[L:E]=\ell$ in weight two. If $r$ is the order of the finite component $\eta$, then $r$ is even and $\varphi(r)\mid 2\ell$ because $\mu_r\subset L$. For even $r$ the possibilities are $r=2,4,6$, and $r=2(2\ell+1)$ if $2\ell+1$ is prime. Indeed $\varphi(r)=2\ell$ requires an odd prime factor of $r$; its predecessor divides $2\ell$ and is even. The possibilities for that prime are $3$ and $2\ell+1$. A power of $3$ cannot supply $\ell\ge7$ in the totient, which gives the assertion.

In the last case $q=2\ell+1$ and $L=\mathbb Q(\zeta_q)$, so its unique quadratic subfield is $E=\mathbb Q(\sqrt{-q})$. This field cannot have class number one: none of the prime absolute discriminants $3,7,11,19,43,67,163$ in the class-number-one list has $(q-1)/2$ prime and at least seven. It cannot have class number $\ell$ either. For any negative prime discriminant $-q$ with $q\ge7$, a reduced form $(a,b,c)$ satisfies $a\le\sqrt{q/3}$; there are at most $a$ odd values $-a<b\le a$, by the boundary convention, and $b$ determines $c$. Hence
\[
             h(-q)\le\sum_{1\le a\le\sqrt{q/3}}a
                     <q/3<(q-1)/2.
\]
The same estimate will also exclude $q=\ell$ when $h(-q)=\ell$. If $r=4$ or $6$, then $E$ must be the quadratic cyclotomic field $\mathbb Q(\mu_r)$, since adjoining any other quadratic field to $E$ would give an even subdegree in $L/E$. Both choices have class number one. In class number one, the values on principal ideals already give $L=E(\mu_r)$, a contradiction here. Finally $r=2$ also precludes $h_E=1$, since then $L=E$. Thus $r=2$ and $h_E=\ell$. Genus theory now gives $D_E=-q$ for a single odd prime $q\equiv3\pmod4$ (the single dyadic possibilities have class number one).

\paragraph{Construction, uniqueness, and distinction.}
Choose a nonprincipal degree-one prime ideal $P$ in $E$ away from $q\ell$, with norm $p$, and write $P^\ell=(\alpha)$. Such primes exist by the prime ideal theorem in ideal classes. Define $D_0(X,p)=2$, $D_1(X,p)=X$, $D_n=XD_{n-1}-pD_{n-2}$ for $n\ge2$. A defining polynomial is
\[
                     D_\ell(X,p)-\operatorname{Tr}_{E/\mathbb Q}\alpha.
\]
Indeed, take the Legendre character on $(\mathcal O_E/(\sqrt{-q}))^\times=\mathbb F_q^\times$ as $\eta_0$. Its rational restriction is $\chi_{-q}$ and its sign on $-1$ ensures $\psi((a))=\eta_0(a)a$ is well defined. Extend to the ideal group by choosing $\psi(P)^\ell=\eta_0(\alpha)\alpha$, using that $P$ generates the cyclic class group. Changing the sign of $\psi(P)$ if necessary, write $\beta^\ell=\alpha$. The polynomial $X^\ell-\alpha$ is irreducible over $E$: its reducibility for prime $\ell$ would make $\alpha$ an $\ell$-th power by the binomial irreducibility criterion, forcing $P$ principal. The extension of $\psi$ therefore has value field $E(\beta)$ of degree $\ell$ over $E$. Conversely, for any admissible $\psi$ of finite component order two, use a representative $xP$ prime to its conductor. The value $\psi(xP)/x$ has $\ell$-th power $\pm\alpha$ and generates $L$ over $E$. Thus all resulting value fields are conjugate. Since $\beta\bar\beta=p$, the real element $t=\beta+p/\beta$ satisfies the stated polynomial. If $t$ were rational, $\beta$ would be at most quadratic over $\mathbb Q$ and $\beta^\ell\in E\setminus\mathbb Q$ would force $\beta\in E$. Hence $t$ generates the real coefficient field of prime degree $\ell$.

To see that different $E$ give different coefficient fields, $q$ can be read off from $K$: it is its unique ramified rational prime away from $\ell$. Indeed $M=E(\zeta_\ell,\beta)$ is Galois over $\mathbb Q$ since it contains the radicals above both $\alpha$ and $\bar\alpha$. Away from $q\ell$, its radical extension over $E(\zeta_\ell)$ is unramified: the valuation of $\alpha$ is a multiple of $\ell$ everywhere, and after removing a local $\ell$-th power one has a unit whose $\ell$-th-root polynomial has unit discriminant. At $q\ne\ell$, inertia maps onto the nontrivial inertia of $E/\mathbb Q$, and fixes $\zeta_\ell$. An inertia lift of conjugation sends $\beta$ to $\zeta_\ell^c p/\beta$ for some $c$. On the $\ell$ distinct conjugates
\[
                 t_j=\zeta_\ell^j\beta+\zeta_\ell^{-j}p/\beta
                 \quad(j\bmod\ell)
\]
it acts by $j\mapsto-j-c$, a nonidentity permutation. Hence $q$ ramifies in the splitting field and in $K$ (a normal closure of a field unramified at $q$ is unramified there). This proves injectivity. \hfill$\square$

The class-number divisibility and the general weight-two finiteness theorem were established by \href{https://arxiv.org/abs/math/0301306}{Yang} and \href{https://ems.press/content/serial-article-files/38349}{Gonz\'alez}. The explicit prime-degree bijection above combines those inputs with the realization and ramification arguments. Its proof is independent of a GRH assumption.

\section{Primitive roots and Euclidean division}

\subsection{Grosswald's conjecture and Erd\H{o}s problem 985}
For a prime $p$, write $g(p)$ for its least positive primitive root and $\widehat g(p)$ for its least prime primitive root. McGown, Trevi\~no and Trudgian establish the following implication in Theorem~1 of their 2016 paper:
\[
 \mathrm{GRH}_{\rm Dirichlet}\ \Longrightarrow\quad
 \begin{cases}
 g(p)<\sqrt{p}-2 & (p>409),\\
 \widehat g(p)<\sqrt{p}-2 & (p>2791).
 \end{cases}
\]
Their paper uses smoothed prime sums to detect a generator of the cyclic group $\mathbb F_p^\times$. For $p\nmid n$ the indicator of a primitive root is
\[
 \frac{\varphi(p-1)}{p-1}
 \sum_{d\mid p-1}\frac{\mu(d)}{\varphi(d)}
                 \sum_{\operatorname{ord}(\chi)=d}\chi(n).
\]
Here $\chi$ runs through Dirichlet characters modulo $p$. This formula follows by writing $n=u^k$ for a generator $u$ and applying inclusion--exclusion to $\gcd(k,p-1)=1$. Lemma~1 of the paper bounds the sums of $\Lambda(n)(1-n/x)$ and of $\Lambda(n)\chi(n)(1-n/x)$ for $n<x$. The first uses ordinary RH and the second uses Dirichlet GRH. A prime power generating $\mathbb F_p^\times$ has a prime base that also generates it. The authors' sieve and finite computations give the two cutoffs. Thus the repository's theorem supplies the analytic hypothesis of their theorem, including for the stronger prime-generator bound. The quantifier over characters here is the reason this use of GRH works. See \href{https://doi.org/10.7169/facm/2016.55.2.5}{McGown--Trevi\~no--Trudgian, Theorem~1}.

\paragraph{Corollary.}
For every odd prime $p$ there is a prime $q<p$ that generates $\mathbb F_p^\times$. This is the odd-prime form of \href{https://www.erdosproblems.com/985}{Erd\H{o}s problem~985}. For $p>2791$ take $q=\widehat g(p)$ and use the second inequality. For the remaining 405 odd primes, trial division enumerates $p$ and the distinct prime divisors $\ell$ of $p-1$. Test successive primes $q<p$ until
\[
      q^{(p-1)/\ell}\not\equiv1\pmod p
      \qquad\text{for every prime }\ell\mid p-1.
\]
The order of $q$ divides $p-1$ and any proper divisor divides $(p-1)/\ell$ for some such $\ell$, proving the test. The accompanying \texttt{checks/primitive\_roots.py} runs this test and independently enumerates all powers of each witness. Every odd prime in the range passes. The largest least witness is 53, attained only at 2791. The literal all-primes formulation has the elementary exception $p=2$. The strict endpoints of the stronger estimates also matter: $g(409)=21$ and $\widehat g(2791)=53$, so the respective estimates fail at their endpoints.

These deductions combine the repository with published implications and a small finite computation. For context, a 2026 large-prime result established the Grosswald bound for $p>1.77\times10^{54}$ without GRH. See \href{https://arxiv.org/abs/2608.20431}{Bagger, \emph{A note on least primitive roots}}.

\subsection{Which cyclic fields admit norm-Euclidean division?}
The ring of integers of $K$ is norm-Euclidean if, given $a,b\in\mathcal O_K$ with $b\ne0$, some $c\in\mathcal O_K$ satisfies $|N_{K/\mathbb Q}(a-bc)|<|N_{K/\mathbb Q}(b)|$. Lezowski and McGown classified these fields in two prime degrees assuming GRH. Their classification yields the following table upon supplying the repository's theorem.
\begin{center}
\begin{tabular}{c|l}
Degree of cyclic $K/\mathbb Q$ & Discriminants of norm-Euclidean $K$\\
\hline
$5$ & $11^4,\ 31^4,\ 41^4$\\
$7$ & $29^6,\ 43^6$\\
$19,31,37,43,47,59,67,71,73,79,97$ & none
\end{tabular}
\end{center}

\paragraph{Why the available GRH suffices.}
For a cyclic field $K/\mathbb Q$ of prime degree $\ell$ and conductor $f$, let $\chi$ be a generator of its group of Dirichlet characters. The abelian factorization is
\[
     \zeta_K(s)=\zeta(s)\prod_{j=1}^{\ell-1} L(s,(\chi^j)^*),
\]
where the star takes the associated primitive character. This convention includes the correct Euler factors at ramified primes. In the open critical strip the factors are holomorphic, so zero-freeness away from the critical line passes to their product. McGown explicitly specifies Dirichlet GRH as the default hypothesis of his conductor estimates, with individual $L(s,\chi)$ or $\zeta_K(s)$ indicated when needed. All these functions fall under the displayed factorization or the repository theorem directly. See \href{https://arxiv.org/abs/1102.2043}{McGown, \emph{Norm-Euclidean Galois fields and the Generalized Riemann Hypothesis}}, introduction and Theorems~3.1--3.2.

\paragraph{Where the classification is proved.}
Proposition~5.2 of Lezowski--McGown uses GRH bounds for character nonresidues and an auxiliary integer to bound the conductor of a norm-Euclidean field. Proposition~2.1 and Table~2.2 give $f<6\times10^{10}$ in degree five and $f<4\times10^{10}$ in degree seven. Proposition~2.2 supplies the finite exclusions, and Proposition~2.3 supplies the Euclidean minima of the surviving fields in these two degrees. Their computations and algebraic arguments give Theorems~1.3 and~1.4. The authors already proved the entire conditional implication; the contribution here is identifying precisely why the formal result meets its input. The extensive field computations are those of the published source. See \href{https://doi.org/10.1090/mcom/3169}{Lezowski--McGown, \emph{The Euclidean algorithm in quintic and septic cyclic fields}}.

Degree eleven illustrates the remaining arithmetic problem. Their Table~1.1 leaves conductors $23,67,331$ as candidates. Theorem~1.5 makes that list exhaustive under the same GRH. It supplies no positive norm-Euclidean certificate for an individual candidate. In degree three the complete classification already has a separate unconditional argument. See \href{https://arxiv.org/abs/2507.05862}{Bagger, Booker, Kerr, McGown, Starichkova and Trudgian}.

\section{A simultaneous weighted prime race}\label{sec:weightedrace}

Chebyshev's original exponential smoothing asks whether
\[
   \sum_{p>2}(-1)^{(p-1)/2} e^{-p/x}\;\longrightarrow\;-\infty .
\]
Hardy--Littlewood and Landau identified this assertion with the zero
location for the nonprincipal character modulo four. The repository's theorem supplies
that input. For the precise equivalence and its history see M.~Suzuki,
\emph{On variants of Chebyshev's conjecture}, \S1.2, equations (16)--(17),
\href{https://doi.org/10.1007/s11139-025-01238-9}{doi:10.1007/s11139-025-01238-9}.
With a different smoothing we can go beyond this single race and
measure the lead in every nonsquare class. The following deduction keeps track of
possible zeros at the center.

\paragraph{Theorem.}
Fix a modulus $q$, a quotient $G$ of $(\mathbb Z/q\mathbb Z)^\times$,
and write $h=|G|$, $G^2=\{a^2:a\in G\}$, $G[2]=\{a:a^2=1\}$.
For $g\in G$ put
\[
 S_g(x)=\sum_{\substack{p\le x,\;p\nmid q\\ {[p]}=g}}
            \frac{\log p}{\sqrt p}\log\frac{x}{p}.
\]
For every nonprincipal $\chi\in\widehat G$ let
$m_\chi=\operatorname{ord}_{s=1/2} L(s,\chi)$, using the character
pulled back to modulus $q$. Then, simultaneously for all $g$,
\[
 S_g(x)-S_1(x)=B_g(\log x)^2+O_q(\log x),
\]
where
\[
 B_g=\frac{|G[2]|}{4h}\,\mathbf 1_{g\notin G^2}
    +\frac1{2h}\sum_{\chi\ne1}
       (1-\operatorname{Re}\chi(g))m_\chi .
\]
In particular, every nonsquare coset eventually leads the identity coset,
by at least $\bigl(|G[2]|/(4h)+o(1)\bigr)(\log x)^2$.
If every element of $G$ has order dividing two, all $g\ne1$ lead by
at least $(1/4+o(1))(\log x)^2$. For a multiquadratic extension of
$\mathbb Q$, the identity class describes primes that split completely,
and the other cosets specify patterns of quadratic splitting at unramified primes.

\paragraph{Proof.}
For a nonprincipal character let $u=\log x$. Suzuki's explicit formulas
(62) and (64), in the proof of Theorem~8 of the cited paper, give
\[
 A_\chi(x)=\sum_{n\le x}\frac{\Lambda(n)\chi(n)}{\sqrt n}\log(x/n)
      =-\frac{m_\chi}{2}u^2+O_q(u).
\]
Indeed the central zero gives the quadratic term. The remaining nontrivial
zeros $1/2+i\gamma$ contribute a bounded series, since
$|x^{i\gamma}|=1$ by Dirichlet GRH and
$\sum_{\gamma\ne0}|\gamma|^{-2}<\infty$ for each fixed character,
with multiplicities. All other terms are $O_q(u)$.
Passing from a primitive character to a character it induces costs
$O_q(u)$ in $A_\chi$, by summing the geometric tails at the finitely many
removed primes. The extra Euler factors do not vanish at $s=1/2$.

The terms $n=p^k$, $k\ge3$, contribute $O(u)$ because
$\sum_p (\log p)p^{-3/2}/(1-p^{-1/2})$ converges. The terms $n=p^2$ equal
\[
 2\int_1^{\sqrt x}
        \left(\sum_{p\le t}\frac{\chi^2(p)\log p}{p}\right)\frac{dt}{t}
 =\mathbf 1_{\chi^2=1}\frac{u^2}{4}+O_q(u).
\]
For principal $\chi^2$ the inner sum is $\log t+O_q(1)$ by Mertens'
estimate. For nonprincipal $\chi^2$ it is $O_q(1)$ by partial summation
of the prime number theorem in progressions at fixed modulus, in its
classical form with error $O_q(t\exp(-c_q\sqrt{\log t}))$.
Subtracting the prime powers thus gives a prime-only sum with quadratic
coefficient $-m_\chi/2-\mathbf 1_{\chi^2=1}/4$.

Character orthogonality expresses $S_g-S_1$ as the sum of those
prime-only sums with coefficients $(\overline{\chi(g)}-1)/h$. The principal
character cancels. Conjugate characters have equal central orders, so the
final coefficient may be written using real parts. Characters satisfying
$\chi^2=1$ are the dual of $G/G^2$. Their sum at $g$ equals
$|G/G^2|$ for $g\in G^2$ and zero otherwise. Since
$|G/G^2|=|G[2]|$, this gives the stated coefficient. Each central-zero
term is nonnegative. The group is finite, so the errors can be taken
simultaneously for all classes.\hfill$\square$

\paragraph{Example modulo eight.}
For each $g=3,5,7$,
\[
  S_g(x)-S_1(x)=\tfrac14(\log x)^2+O(\log x).
\]
Here all three central orders are zero, as one can see without a
numerical calculation. On the odd classes $1,3,5,7$ the three respective
sign patterns are $+,-,+,-$, $+,+,-,-$, and $+,-,-,+$.
The Dirichlet series at $s=1/2$ converges by periodicity and zero mean.
Each consecutive block of four odd summands has positive sum for all
three patterns. For the last pattern this uses
$f(a)-f(a+2)>f(a+4)-f(a+6)$ with $a=8k+1$ and
$f(t)=t^{-1/2}$, which follows from strict convexity. Thus every central
value is positive. The example gives three strict eventual inequalities
with a common leading coefficient. The theorem concerns the specified
weighted sums; an unweighted two-way prime race can change sign.

\paragraph{Provenance.}
Suzuki's Theorems~3 and~4 already give modulo-four conclusions and his
Theorem~9 gives the individual real-character prime sum for a nonzero
central value. The quotient formula above combines his explicit formula
with the square map on $G$; it is presented as a deduction without a claim
of priority.

\section{A one-sided test for a nonabelian cubic field}\label{sec:cubic}
Let $E=\mathbb Q(\sqrt[3]2)$ and $M=E(\zeta_3)$. Our target is RH for
$\zeta_E$, already beyond the public theorem. The polynomial $X^3-2$
is Eisenstein at two and has discriminant $-108$. Its three roots give
$\operatorname{Gal}(M/\mathbb Q)=S_3$, with quadratic subfield
$\mathbb Q(\sqrt{-3})$. For $p>3$ write $g_p$ for its Frobenius conjugacy
class, and $r_p$ for the number of roots modulo $p$. The following table
records all three possibilities. Write $\tau$ for the character of the
two-dimensional standard representation.
\begin{center}
\begin{tabular}{lrrrr}
\toprule
Factorization mod $p$ & Class size & $r_p$ & $\operatorname{sgn}(g_p)$ & $\tau(g_p)$\\
\midrule
Three linear factors & 1 & 3 & $1$ & $2$\\
One linear and one quadratic factor & 3 & 1 & $-1$ & $0$\\
One irreducible cubic & 2 & 0 & $1$ & $-1$\\
\bottomrule
\end{tabular}
\end{center}
For example, $7$ and $31$ are both $7\pmod{24}$. The cubes modulo seven are $0,1,-1$,
whereas $4^3\equiv2\pmod{31}$ and all three cube roots of unity lie in $\mathbb F_{31}$.
The two traces are therefore $-1$ and $2$. In general $\tau(g_p)=r_p-1$. The character of the permutation action
on the roots is $1+\tau$. Its Artin $L$-function is $\zeta_E$, so the
standard factor is $F=\zeta_E/\zeta$. In the right half of the critical
strip $F$ is holomorphic using the repository's ordinary RH, and
$F(1)$ is finite and positive by the simple poles of both zeta functions.
This local argument suffices everywhere we use $F$.

\paragraph{Theorem (three splitting types).}
Use Dirichlet GRH as established in the pinned source. Choose any real weights
$w=(u,v,z)$ on the three rows, in order. Put
\[
 A=(u+3v+2z)/6,\quad B=(u-3v+2z)/6,\quad C=(u-z)/3.
\]
Thus $w=A+B\operatorname{sgn}+C\tau$. Set
\[
 S_w(x)=\sum_{3<p\le x}w(g_p)\log p\,(1/p-1/x),\qquad
 R_w(x)=S_w(x)-A\log x-c_w.
\]
In this sum and below $p$ ranges over primes. The constant $c_w$ is fixed by
$R_w(x)\to0$. The following statements classify the tests obtained from these weights.
\begin{enumerate}
\item If $C=0$, then $R_w(x)=O_w(x^{-1/2}\log^2(2x))$ follows from Dirichlet GRH.
\item If $C\ne0$, RH for $\zeta_E$ is equivalent to the following
\emph{one-sided} estimate. For every $\epsilon>0$ there exists
$K_{w,\epsilon}>0$ such that, for all sufficiently large $x$,
\[
 \operatorname{sgn}(C) R_w(x)\le K_{w,\epsilon}x^{-1/2+\epsilon}.
\]
Under RH for $\zeta_E$ one also gets the two-sided error in part 1 for these weights.
\end{enumerate}
For example, just an upper bound for primes at which $X^3-2$ splits
completely suffices. Just a lower bound for primes at which it is
irreducible also suffices. For $w=(2,0,-1)$ the logarithmic main term
vanishes. This last choice measures exactly the part missed by the two
Dirichlet characters in the table.

\paragraph{Constants and direct estimates.}
Let $\theta_w(t)=\sum_{3<p\le t}w(g_p)\log p$. Integration gives
$S_w(x)=\int_1^x\theta_w(t)t^{-2}\,dt$. The fixed-field prime ideal
theorem for $E$, $\mathbb Q(\sqrt{-3})$ and $\mathbb Q$ gives
$\theta_w(t)=At+O_w(t e^{-b\sqrt{\log t}})$ for some $b>0$ and large $t$.
Indeed degree-one primes of $E$ above $p>3$ are counted by $r_p$, and the
quadratic field similarly gives $1+\operatorname{sgn}(g_p)$.
Terms of higher residue degree and higher ideal powers add
$O(\sqrt t\log^2(2t))$ in each fixed field, by summing over $p^j\le t$,
$j\ge2$. A possible fixed exceptional zero contributes $O(t^{\beta_0})$,
which is absorbed for large $t$. Thus
\[
 c_w=\int_1^\infty(\theta_w(t)-At)t^{-2}\,dt,\qquad
 R_w(x)=-\int_x^\infty(\theta_w(t)-At)t^{-2}\,dt .
\]
Under RH for the respective Dedekind zeta function, the prime ideal
theorem has error $O(\sqrt t\log^2(2t))$. For the quadratic field this
follows from $\zeta_{\mathbb Q(\sqrt{-3})}=\zeta L(s,\chi_{-3})$.
The last integral now proves part 1, and the forward assertion in part 2.
For the prime ideal estimate see
\href{https://arxiv.org/abs/1312.4463}{Greni\'e and Molteni, \emph{Explicit versions of the prime ideal theorem}}.

\paragraph{Why the reverse implication needs only one side.}
By linearity and part 1 it suffices to work with $w=\tau$ and an upper
bound. Initially for $\Re s>1$ set
\[
 D(s)=\sum_{p>3}(r_p-1)(\log p)p^{-s}.
\]
The Euler factors give $D(s)=-F'(s)/F(s)-H(s)$ with $H$ holomorphic
on $\Re s>1/2$. Indeed the $j$th power coefficient of an unramified
prime is $\tau(g_p^j)\log p$ and $|\tau|\le2$. The sum for $j\ge2$
converges absolutely locally uniformly there. The two ramified primes
have locally holomorphic corrections, since their local inverse roots are
zero or roots of unity. The prime ideal theorem above gives
$c_\tau=D(1)$ by Abelian passage to the limit, or by integrating
$D(s)=s\int_1^\infty\theta_\tau(x)x^{-s-1}\,dx$ down to $s=1$.

Use the repository's smoothing and take a Mellin transform. Absolute convergence for
$\Re s>1$ gives, by integrating from $p$ to infinity,
\[
 T(s)=\int_1^\infty R_\tau(x)x^{-s}\,dx
      =\frac{D(s)}{s(s-1)}-\frac{D(1)}{s-1}.
\]
The apparent singularity at one is removable. Each zero $\rho$ of $F$ in
$\Re s>1/2$ supplies a simple pole of $T$ with residue
$m_\rho/[\rho(1-\rho)]$, where $m_\rho$ is the multiplicity. In particular
every real pole in that half-plane has \emph{positive} residue, since real
zeros there lie strictly below one.

Fix $1/2<\beta<1$ and take the assumed bound with $\epsilon=\beta-1/2$.
The function $U(x)=K x^{\beta-1}-R_\tau(x)$ is nonnegative beyond a finite
$X$. Its Mellin integral from $X$ to infinity equals
$K/(s-\beta)-T(s)$ up to an entire function. Suppose its abscissa of
convergence $\sigma_c$ exceeds $\beta$. It is at most one since
$R_\tau(x)\to0$. Landau's lemma for a nonnegative function forces a
singularity at the real point $\sigma_c$. Meromorphic continuation here
permits only a simple pole with \emph{negative} residue, and this would make
the integral negative as $s\downarrow\sigma_c$ on the real axis. This
contradiction shows $\sigma_c\le\beta$. The integral is holomorphic
throughout $\Re s>\beta$, so $T$ has no poles there. Taking all $\beta>1/2$
proves that $F$ is zero-free on the right, and hence so is $\zeta_E$ away
from its pole. The Dedekind functional equation gives the left half.
This proves the equivalence. In this argument the prime-power correction
is holomorphic on the \emph{open} half-plane, which is precisely the
domain needed for the conclusion.

\paragraph{An exact limit on what congruence weights detect.}
Even a very large fixed modulus cannot recover the missing trace.
For any complex function $a$ on the units modulo a fixed $q$,
Chebotarev gives
\[
 \lim_{x\to\infty}\frac1{\pi(x)}
 \sum_{\substack{p\le x\\p\nmid 6q}} |\tau(g_p)-a(p\bmod q)|^2
   =1+\frac{1}{\varphi(q)}\sum_{b\in(\mathbb Z/q)^\times}|a(b)|^2 .
\]
To prove this, form $M\mathbb Q(\zeta_q)$. The intersection of the
two fields is abelian over $\mathbb Q$, so it is either $\mathbb Q$ or the
quadratic subfield of $M$. Conditional on a cyclotomic Frobenius, the
$S_3$ component is uniform on the entire group or on a coset of $A_3$.
The trace has mean zero on the even coset ($2,-1,-1$) and on the odd
coset ($0,0,0$). Its unconditional mean square is $(4+2)/6=1$.
Expanding the square proves the formula, because Chebotarev equidistributes
the compatible pairs with their conjugacy multiplicities.
In particular, linear combinations of any finite set of Dirichlet characters
cannot distinguish this trace from zero in their correlations.
This rules out an exact encoding by fixed residue classes. Weighted
estimates exploiting genuinely new arithmetic structure remain possible.

\paragraph{What this would give for Artin's conjecture.}
For base $2$, Hooley's argument uses the fields
$\mathbb Q(\zeta_n,2^{1/n})$ for squarefree $n$. Our $M$ is the instance
$n=3$. Its zeta function factors as
$\zeta_M=\zeta\,L(s,\chi_{-3})F^2$ by the regular representation of $S_3$.
Thus the criterion identifies one concrete missing instance, including
its precise sign and arithmetic weight. Resolving the full primitive-root
problem by this route would require the family and the analytic control
used in its sieve. For the precise conditional theorem see
\href{https://arxiv.org/abs/2505.05601}{Fan and Pollack, \emph{Counting primes with a given primitive root, uniformly}}.

\section{Further research questions}
\begin{enumerate}[leftmargin=*]
\item In coefficient degree four, Yang's divisibility allows CM fields of
class number one, two or four. How do the quartic coefficient fields from
the cyclic and noncyclic class groups of order four overlap? The odd-degree
proof uses that a sign has an odd root and that the coefficient field has
a real embedding; neither argument covers every quartic field.
\item For the quintic fields, determine the exact conductor and discriminant
at five uniformly from the local fifth-power class of the generator. This
would complement the defining-polynomial table with an intrinsic discriminant formula.
\item Can the upper/lower covolume comparison produce a bound for the
weight $r_p-1$ of Section~\ref{sec:cubic}? One would need a local
divisibility formula sensitive to the difference between the identity and a
three-cycle, and a matching archimedean main term and denominator estimate.
The orthogonality calculation there rules out expressing this weight by a
fixed family of residue classes alone.
\end{enumerate}
\end{document}
